% Compile with pdfLaTeX. No external assets or custom fonts required.
% Content tailored only from the supplied CV and stated graduation term.
\documentclass[11pt,letterpaper]{article}
\usepackage[margin=0.60in,top=0.53in,bottom=0.55in]{geometry}
\usepackage[T1]{fontenc}
\usepackage[utf8]{inputenc}
\usepackage{mathptmx}
\usepackage{titlesec,enumitem,xcolor,hyperref,fancyhdr}
\input{glyphtounicode}
\pdfgentounicode=1
\definecolor{navy}{RGB}{29,52,75}
\hypersetup{colorlinks=true,urlcolor=navy,linkcolor=navy,pdfauthor={Wanli Qian},pdftitle={Wanli Qian -- Summer 2027 Research Internship}}
\pagestyle{fancy}
\fancyhf{}
\renewcommand{\headrulewidth}{0pt}
\fancyfoot[L]{\footnotesize\color{navy}Wanli Qian}
\fancyfoot[R]{\footnotesize\thepage\ / 2}
\setlength{\footskip}{18pt}
\setlength{\parindent}{0pt}
\setlength{\parskip}{0pt}
\setlength{\emergencystretch}{2em}
\titleformat{\section}{\normalsize\bfseries\color{navy}}{}{0em}{}[\vspace{-4pt}\rule{\textwidth}{0.45pt}]
\titlespacing*{\section}{0pt}{9pt}{4pt}
\setlist[itemize]{leftmargin=13pt,label=\textbullet,itemsep=2pt,topsep=3pt,parsep=0pt,partopsep=0pt}
\newcommand{\role}[2]{\textbf{#1}\hfill\textit{#2}\par}
\newcommand{\pub}[3]{\textbf{#1}\par #2\par #3\par\vspace{7pt}}
\begin{document}
\fontsize{10.5}{12.4}\selectfont
{\fontsize{24}{27}\selectfont\textbf{Wanli Qian}}\par
\vspace{3pt}
\href{mailto:wanliqia@usc.edu}{wanliqia@usc.edu} \enspace|\enspace (678) 956-4378 \enspace|\enspace Los Angeles, CA\par
\href{https://silverwings-zero.github.io/}{Website} \enspace|\enspace
\href{https://www.linkedin.com/in/wanli-qian-7bb399192/}{LinkedIn} \enspace|\enspace
\href{https://github.com/silverwings-zero}{GitHub} \enspace|\enspace
\href{https://scholar.google.com/citations?user=azFoAIkAAAAJ&hl=en}{Google Scholar}\par
\vspace{4pt}
\textbf{Seeking Summer 2027 Research Internship}\par
\textcolor{navy}{Generative AI | Multimodal Learning | Physical Interaction}\par
\section*{Research Profile}
CS Ph.D. researcher developing language-conditioned generative models and interactive systems that connect language, vision, and touch. Builds multimodal haptic datasets and AI-driven authoring tools for physical and virtual environments. First-authored work includes IEEE Haptics Symposium 2026 (paper award) and co-first-authored ACM UIST 2024 research.\par

\section*{Education}
\role{University of Southern California}{2024--Present}
Ph.D. in Computer Science \hfill \textbf{Expected graduation: Fall 2028}\par
\vspace{3pt}
\role{University of Chicago}{2022--2024}
Predoctoral M.S. in Computer Science\par
\vspace{3pt}
\role{Georgia Institute of Technology}{2018--2022}
B.S. in Computer Science\par
\section*{Research Experience}
\role{USC HaRVI Lab | Ph.D. Researcher}{2024--Present}
\begin{itemize}
\item Develop language-conditioned generative models that synthesize tactile and vibration signals from force, acceleration, and vision-based tactile data.
\item Curate multimodal haptic datasets spanning 200+ materials, combining tapping, sliding, audio, and affective descriptors to support model learning.
\item Investigate how robots acquire, represent, and generate haptic understanding through multimodal sensing, language, and action.
\item Integrate VR/AR and language-based interaction for haptic and texture authoring; build tactile-feedback teleoperation systems for perception and data collection.
\item Study exploratory robot policies that adapt actions to reduce uncertainty about object and surface properties.
\end{itemize}
\vspace{3pt}
\role{University of Chicago AxLab | Predoctoral Researcher}{2022--2024}
\begin{itemize}
\item Led SHAPE-IT development for on-the-fly 3D shape generation and physical reconfiguration from conversational language input.
\item Developed adaptive LLM-based agents that modify control scripts during user interaction; investigated AI-driven authoring for tangible and shape-changing interfaces.
\end{itemize}
\vspace{3pt}
\role{Georgia Tech Borg Lab | Lab Assistant}{2021--2022}
\begin{itemize}
\item Explored computational representations of artistic strokes, including polynomial models and deformable virtual brush dynamics.
\item Combined motion capture, trajectory optimization, and robot control to reproduce human drawing and writing motions with physical robots.
\end{itemize}
\section*{Industry Experience}
\role{Kolmostar | Software Engineer Intern}{2022}
\begin{itemize}
\item Built simulation pipelines modeling urban geometry and wireless-signal attenuation; created interactive visualization tools for analysis and deployment planning.
\end{itemize}
\vspace{3pt}
\role{Megvii | Robotics Software Solution Intern}{2019}
\begin{itemize}
\item Implemented and benchmarked probabilistic and graph-based motion planners for industrial manipulators, with visualizations of planner behavior across task geometries.
\end{itemize}
\newpage
\section*{Selected Publications}
\pub{Language-Guided Multimodal Texture Authoring via Generative Models.}{\textbf{Wanli Qian}, Michael Gu, Aiden Chang, Shihan Lu, Heather Culbertson.}{\textit{IEEE Haptics Symposium, 2026.} \enspace \href{https://arxiv.org/abs/2604.06489}{Paper} \enspace|\enspace \href{https://youtu.be/f5vz52uPw-4}{Video}\par\textbf{Distinguished Technical Long Paper Award}}
\pub{Shape n'Swarm: Hands-on, Shape-aware Generative Authoring with Swarm UI and LLMs.}{Matthew Jeung, Anup Sathya, \textbf{Wanli Qian}, Steven Arellano, Luke Jimenez, Ken Nakagaki.}{\textit{ACM UIST, 2025.} \enspace \href{https://dl.acm.org/doi/10.1145/3746059.3747781}{Paper} \enspace|\enspace \href{https://www.youtube.com/watch?v=5u0M9yL7tyY}{Video}}
\pub{SHAPE-IT: Exploring Text-to-Shape-Display for Generative Shape-Changing Behaviors with LLMs.}{\textbf{Wanli Qian}*, Chenfeng Gao*, Anup Sathya, Ryo Suzuki, Ken Nakagaki.}{\textit{ACM UIST, 2024. *Equal contribution.} \enspace \href{https://dl.acm.org/doi/10.1145/3654777.3676348}{Paper} \enspace|\enspace \href{https://www.youtube.com/watch?v=2NxzOBc7AdQ}{Video}}
\pub{Weakly Supervised Part-Based Method for Combined Object Detection in Remote Sensing Imagery.}{\textbf{Wanli Qian}, Zhe Yan, Zhe Zhu, Wei Yin.}{\textit{IEEE JSTARS, vol. 15, pp. 5024--5036, 2022.} \enspace \href{https://ieeexplore.ieee.org/document/9789416}{Paper}}
\pub{Dynamic Pseudo-Label Generation for Weakly Supervised Object Detection in Remote Sensing Images.}{Hui Wang, Hao Li, \textbf{Wanli Qian}, Wenhui Diao, Liangjin Zhao, Jinghua Zhang, Daobing Zhang.}{\textit{Remote Sensing, vol. 13, no. 8, p. 1461, 2021.} \enspace \href{https://www.mdpi.com/2072-4292/13/8/1461}{Paper}}
\section*{Manuscripts Under Review}
\pub{TouchTwin: Human-in-the-Loop Haptic Texturing of 3D-Scanned Tabletop Scenes with Vision--Language Recognition and Language-Guided Generation.}{\textbf{Wanli Qian} et al.}{\textit{Submitted to ACM CHI 2027; under review.}}
\pub{Feeling What the Gripper Feels: Haptic Rendering of a Grasped Object from Tactile Video in Teleoperation.}{\textbf{Wanli Qian} et al.}{\textit{Submitted to IEEE Robotics and Automation Letters (RA-L); under review.}}

\section*{Teaching}
\textbf{CSCI 445: Introduction to Robotics} | Teaching Assistant, USC \hfill 2025\par
\vspace{3pt}
\textbf{CMSC 25040: Introduction to Computer Vision} | Teaching Assistant, UChicago \hfill 2023\par
\section*{Service \& Outreach}
\textbf{Reviewing:} CHI, UIST, IEEE Haptics Symposium, IEEE Transactions on Haptics.\par
\vspace{3pt}
\textbf{Mentorship:} USC CURVE (Spring/Fall 2025, Spring/Fall 2026); USC Viterbi SURE (Summer 2025); USC REU (Summer 2026).\par
\vspace{3pt}
\textbf{Outreach:} Peace4Kids robotics workshop organizer/mentor; RSS conference volunteer.\par
\end{document}
